% THE MAP. One three-column geometry, drawn by \rulemap, appears on slides 3, 15 and 26.
% The same \providecommand block opens 03, 15 and 26; the first one loaded defines it, so all three
% slides draw identical headers, arrows, sub-headers and body anchors. Edit all three together.
% Arguments: #1 #2 #3 = column bodies (compute once, count once, publish once); #4 = extra TikZ.
\providecommand{\rulemap}[4]{%
\noindent\begin{tikzpicture}[x=1cm,y=1cm,
  rh/.style={evidence,text width=4.2cm,minimum height=.9cm,font=\fontsize{12}{14}\selectfont\bfseries,text=Teal},
  rs/.style={anchor=north,inner sep=0pt,font=\fontsize{10}{12}\selectfont,text=Muted},
  rb/.style={anchor=north,inner sep=0pt,text width=4.7cm,align=center,font=\fontsize{12}{14.5}\selectfont}]
\useasboundingbox (0,1.05) rectangle (14.4,-3.75);
\node[rh] (r1) at (2.3,0) {1\enspace COMPUTE ONCE};
\node[rh] (r2) at (7.2,0) {2\enspace COUNT ONCE};
\node[rh] (r3) at (12.1,0) {3\enspace PUBLISH ONCE};
\draw[flow] (r1)--(r2);\draw[flow] (r2)--(r3);
\draw[message] (r3.north)--++(0,.42)-|node[pos=.25,above,inner sep=1.5pt,font=\fontsize{9}{10}\selectfont,text=Teal]{continue} (r1.north);
\node[rs] at (2.3,-.62) {Runtime};
\node[rs] at (7.2,-.62) {Evidence};
\node[rs] at (12.1,-.62) {Acceptance};
\node[rb] at (2.3,-1.15) {#1};
\node[rb] at (7.2,-1.15) {#2};
\node[rb] at (12.1,-1.15) {#3};
#4
\end{tikzpicture}\par}
\begin{frame}[t]{\ev{\tagP}\surf{THREE RULES}Sharing work creates three rules.}
\vspace{.1cm}
\rulemap{When does reusing a state beat rebuilding it?}
{What must each parallel trial return?}
{Which exact bytes may we publish, and who decides?}
{\foreach \xx/\lab in {2.3/a cost rule,7.2/an evidence contract,12.1/an acceptance gate}{\draw[flow,draw=Teal] (\xx,-2.3)--(\xx,-2.68);\node[anchor=north,inner sep=0pt,font=\fontsize{12}{14.5}\selectfont\bfseries,text=Teal] at (\xx,-2.8) {\lab};}}
\claim{\centering\forktree\enspace You will leave with all three.}
\sources{The talk takes the rules in order. This map returns after the examples and at the end.}
\qadetail{This slide restates the canonical three system responsibilities. Runtime starts private continuations and supplies the state the next step reads, priced against a warm alternative such as a build cache. An evaluator or recorder outside child-controlled state binds executed checks to exact artifact digests and records which observations trials share. Acceptance selects or composes candidates after the appropriate check; the controller retains publication authority. Once is the goal of each obligation, not an exactly-once delivery guarantee: compute once means do not repeat preparation when a saved state supplies what the next step reads; count once means a shared observation contributes once to a combined result; publish once means one controller publishes one checked artifact and rejects stale results (the epoch fence is proposed). The contribution claim is the integrated execution, evidence and acceptance contract; forking, testing, review and provenance are established techniques, and no claim of invention or measured parallel speedup is made.}
\note{[0:50] Those obligations are three rules, and they organize this talk. Compute once: reuse a state only when that beats rebuilding it, even from a warm cache. Count once: parallel trials share inputs and errors, so a shared observation counts once. Publish once: one controller publishes the exact bytes that passed the check. Each rule answers one question, and you will leave with a cost rule, an evidence contract and an acceptance gate. After five examples, we take the rules in order. First, where these rules came from.}
\end{frame}
